1.4. 设 $V$ 为秩为 $2m>0$ 的局部自由 $A$-模, $Q$ 为 $V$ 上的非退化二次型. 假定给定一个分解 $V=W_1\oplus W_2$, 其中两子空间均为全奇异的(totally singular). 此时型 $Q(w_1+w_2)$ 使 $W_1$ 和 $W_2$ 互为完美对偶. 回顾对于 $f\in W_2$ (等同于 $W_1^*$), 与 $f$ 的缩并是 $\bigwedge W_1$ 上由下式定义的 $-1$ 次自同态

\begin{equation*}
f\mathbin{\lrcorner}(e_1\wedge\cdots\wedge e_p)
=\sum_{i=1}^{p}(-1)^{i-1}\langle e_i,f\rangle
e_1\wedge\cdots\wedge\widehat e_i\wedge\cdots\wedge e_p.
\end{equation*}

\footnote{印本公式中为 $(-1)^i$, 且余下的楔积以未定义的 $e_0$ 开始. 英文版恢复了标准的缩并公式 $(-1)^{i-1}e_1\wedge\cdots\widehat e_i\cdots\wedge e_p$, 紧接其后给出的 Clifford 作用要求此公式. 英文版通过针对第 65 页的高精度局部放大核实了这一读法.}

在 $A$-模 $\bigwedge W_1$ 上赋予如下 $\mathbb Z/2$-分次

%%%%%%%%%% source scan idx 73 | folio 66 %%%%%%%%%%

\begin{equation*}
(\bigwedge W_1)^+=\bigoplus_i\bigwedge^{2i}W_1,
\qquad
(\bigwedge W_1)^-=\bigoplus_i\bigwedge^{2i+1}W_1.
\end{equation*}

于是代数 $\operatorname{End}(\bigwedge W_1)$ 成为 $\mathbb Z/2$-分次代数: 其偶自同态 (分别地, 奇自同态) 是保持分次的自同态 (分别地, 次数模二为一的自同态).

Bourbaki, \emph{Algèbre}, 第~9 章, \S9, 第~4 节证明了存在唯一的 $\mathbb Z/2$-分次代数同构 $s$, 从 $C(Q)$ 到 $\operatorname{End}(\bigwedge W_1)$, 使得

\begin{equation*}
s\bigl(\rho((e,0))\bigr)=e\wedge(-),
\qquad
s\bigl(\rho((0,f))\bigr)=f\mathbin{\lrcorner}(-).
\end{equation*}

暂以 $e(W_1,W_2)$ 表示 $C^+(Q)$ 中使 $s(e(W_1,W_2))$ 在 $(\bigwedge W_1)^+$ 上为恒等映射、在 $(\bigwedge W_1)^-$ 上为零的元素. $C^+(Q)$ 的中心 $Z(Q)$ 同构于 $A\oplus A$, 这一分解由幂等元 $e(W_1,W_2)$ 和 $1-e(W_2,W_1)$ 定义. 此外, $C^+(Q)$ 是 $Z(Q)$ 上的矩阵代数.

由此和 1.2, 如同上述出处中一样可推出:

\underline{命题} 1.5. \underline{在 1.1 的假设下, 若 $n=2m>0$ 且 $Q$ 非退化, 则 $C^+(Q)$ 的中心 $Z(Q)$ 是 $A$ 上秩为二的局部自由平展代数, 而 $C^+(Q)$ 是 $Z(Q)$ 上的 Azumaya 代数.}

1.6. 设 $V$ 和 $Q$ 如 1.4, $W$ 为 $V$ 的全奇异直和项, 且局部自由、秩为 $m$. 设 $W'_1$ 为 $W$ 的一个补, 因而 $V=W\oplus W'_1$. $W$ 的每个补 $W'_2$ 均可等同于映射 $u:W'_1\to W$ 的图. 型 $\Phi$ 使 $W$ 与 $W'_1$ 互为对偶, 故 $u$ 由 $W'_1$ 上的双线性型

%%%%%%%%%% source scan idx 74 | folio 67 %%%%%%%%%%

\begin{equation*}
B(x,y)=\Phi(x,u(y)),
\end{equation*}

唯一确定, 且 $W'_2$ 全奇异当且仅当

\begin{equation*}
B(x,x)=-Q(x)
\end{equation*}

在 $W'_1$ 上成立. 由于这样的型 $B$ 总是存在, $W$ 有全奇异的补. 选定其中一个 $W'$ 后, $W$ 的全奇异补 $W'_0$ 的集合可等同于 $W'$ 上交替型的集合, 或等价地, 等同于 $S$ 上仿射空间 $\mathbb E(\bigwedge^2W')$ 的截面集.

对 $\mathbb E(\bigwedge^2W')$ 的每个截面 $x$, 设 $W'_0(x)$ 为 $W$ 的相应补, 并令 $e(x)=e(W,W'_0(x))$. 函数 $x\mapsto e(x)$ 与任意标量扩张相容, 从而定义了 $S$-概形的态射

\begin{equation*}
e:\mathbb E(\bigwedge^2W')\longrightarrow\Spec(Z(Q)).
\end{equation*}

$S$-概形 $\mathbb E(\bigwedge^2W')$ 平坦且纤维几何连通, 而 $\Spec(Z(Q))$ 在 $S$ 上平展. 因而态射 $e$ 经由 $S$ 分解, 从而经由 $\Spec(Z(Q))$ 在 $S$ 上的一个截面分解; 等价地, $e(x)$ 与 $x$ 无关.
\footnote{印本的这句话说 $e$ 经由仿射空间 $\mathbb E(\bigwedge^2W')$ 的一个截面分解. 从连通概形到平展概形的论证实际上给出经由 $S$ 和 $\Spec(Z(Q))$ 的一个截面的分解, 正如下文结论所要求的.}
这就说明下列记号是合理的.

\underline{记号} 1.7. \underline{对于如上的 $W$ 及 $W$ 的任意全奇异补 $W'$, 以 $e(W)$ 表示幂等元 $e(W,W')\in Z(Q)$.}

1.8. 在 1.3 中已见, 当 $\lambda$ 可逆时, $C^+(Q)$ 与 $C^+(\lambda Q)$ 可等同. 与上述类似的连续性论证表明, 此等同将

\begin{equation*}
e(W)\in Z(Q)\qquad\text{映至}\qquad e(W)\in Z(\lambda Q).
\end{equation*}

\footnote{印本中的交叉引用为 1.6; 标量变换的等同是 1.3.1 中的构造.}

%%%%%%%%%% source scan idx 75 | folio 68 %%%%%%%%%%

1.9. 回顾若 $B$ 和 $C$ 为两个 $\mathbb Z/2$-分次代数, 则其分次张量积 $B\otimes' C$ 中的乘法对于齐次的 $b_i$ 和 $c_i$, 由 Koszul 规则

\begin{equation*}
(b_1\otimes c_1)(b_2\otimes c_2)
=(-1)^{\deg(c_1)\deg(b_2)}b_1b_2\otimes c_1c_2
\end{equation*}

定义.

若 $(V_1,Q_1)$ 和 $(V_2,Q_2)$ 为两个二次模, 则

\begin{equation*}
C(Q_1\oplus Q_2)=C(Q_1)\otimes' C(Q_2)
\end{equation*}

(Bourbaki, 同上).

若 $P$ 和 $Q$ 为两个有限型局部自由 $\mathbb Z/2$-分次 $A$-模, 则对齐次的 $u,v,p,q$ 作如下规定, 将 $\mathbb Z/2$-分次代数 $\operatorname{End}(P)\otimes'\operatorname{End}(Q)$ 与 $\operatorname{End}(P\otimes Q)$ 等同:

\begin{equation*}
(u\otimes v)(p\otimes q)
=(-1)^{\deg(v)\deg(p)}u(p)\otimes v(q).
\end{equation*}
